% h2upreamble.tex
\documentclass[lualatex]{jlreq}
% 00preamble.tex
% 基本とドライバ関連
\usepackage{graphicx}
\usepackage[table]{xcolor}
\usepackage{geometry}
\usepackage{siunitx}
\usepackage{multirow}
\usepackage{float}
\usepackage{url}
\usepackage{listings}
\usepackage{caption2}
\usepackage{lscape}
\usepackage{pdflscape}
\usepackage{csvsimple}
\usepackage{wrapfig}
\usepackage{comment}
\usepackage{cprotect}
\usepackage{tikz}
\usetikzlibrary{intersections, calc, arrows.meta}
\usetikzlibrary{positioning, through}
\usepackage{fancybox}
\usepackage{booktabs}
\usepackage{adjustbox}
\usepackage{longtable}
\usepackage{nameref}
% 数式系パッケージ
\usepackage{amsmath,amssymb,amsfonts}
\usepackage{mathtools,amssymb,amsfonts}
\usepackage{amsthm}
\usepackage{bm}
% LuaTeX-ja設定
\usepackage{luatexja}
\usepackage[no-math]{luatexja-fontspec}

\setcounter{section}{0}
\makeatletter
    \renewcommand{\theequation}{
    \thesection.\arabic{equation}}
    \@addtoreset{equation}{section}
    \renewcommand{\figurename}{Figure }
    \renewcommand{\tablename}{Table }
\makeatother
\geometry{top=15truemm,bottom=25.4truemm,inner=19.05truemm,outer=19.05truemm}

% Remove section number from equation number, so that it is just 1, 2, NOT 1.1, 2.1:
% \counterwithout{equation}{section}

% Source code listing settings (listings package)
\lstset{
    frame=tb,
    basicstyle=\ttfamily\small,
    % basicstyle=\ttfamily\small\bfseries,
    columns=fullflexible,
    keepspaces=true,
    numbers=left,
    firstnumber=1,
    numberstyle=\scriptsize,
    breaklines=true,
    % frame=none
    % frame=single,
    % gobble=4,
    % xleftmargin=0.2\textwidth,
    % xrightmargin=0.2\textwidth,
}
% \renewcommand{\lstlistlistingname}{List of Source Codes}
% \renewcommand{\lstlistingname}{リスト}

\NewDocumentCommand{\bissectrice}{%
    O{}     % drawing options
    mmm     % bissector of mmm
    m       % intersection point between base and bissector
    O{1}O{1}% extended drawing of the bissector
    }{%
    \path[name path=Bis#2#3#4] let
        \p1 = ($(#2) - (#3)$),          % chktex 46
        \p2 = ($(#4) - (#3)$),          % chktex 46
        \n1 = {veclen(\x1,\y1)},        % chktex 36
        \n2 = {veclen(\x2,\y2)},        % chktex 36
        \n3 = {max(\n1,\n2)},           % chktex 36
        \p1 = ($(#3)!\n3!(#2)$),        % chktex 46
        \p2 = ($(#3)!\n3!(#4)$),        % chktex 46
        \p3 = ($(\p1) + (\p2) - (#3)$)  % chktex 46
    in
        (#3) --(\p3);                   % chktex 36

    \path[name path = foo] (#2) --(#4); % chktex 36

    \path[name intersections={of=foo and Bis#2#3#4, by=#5}];

    \path[#1] ($(#3)!#6!(#5)$) --($(#5)!#7!(#3)$); % chktex 46, chktex 36
}


\counterwithout{equation}{section}
\geometry{top=20truemm,bottom=25.4truemm,inner=19.05truemm,outer=19.05truemm}
\title{\bfseries {\huge Testing My Preamble for Lua{\LaTeX{}}}\\[0.5em]{\Large Lua{\LaTeX{}}用のプレアンブルのテスト}}
\author{{\LARGE Ivan Setiawan\thanks{MyCompany, Inc.}} \and {\LARGE My Self\,\thanks{株式会社マイカンパニー}}}
\date{Created: 2026年3月28日\\[0.3em]Last Updated: 2026/04/19}

\begin{document}
\maketitle
\hrule
\lstlistoflistings{}
\vspace{0.5em}
\hrule
\vspace{1em}

\section{00preambleの使い方}

Copy and paste the code in Listing~\ref{lst:inputfile} to \texttt{inputfile.tex} and compile it with Lua{\LaTeX}.

\begin{lstlisting}[
 language={[LaTeX] TeX},
 caption={Template \LaTeX{} file (\texttt{inputfile.tex}) to use 00preamble},
 label={lst:inputfile},
 % basicstyle=\ttfamily \footnotesize,
 % firstnumber=3,
]
% inputfile.tex
\documentclass[lualatex]{jlreq}
% 00preamble.tex
% 基本とドライバ関連
\usepackage{graphicx}
\usepackage[table]{xcolor}
\usepackage{geometry}
\usepackage{siunitx}
\usepackage{multirow}
\usepackage{float}
\usepackage{url}
\usepackage{listings}
\usepackage{caption2}
\usepackage{lscape}
\usepackage{pdflscape}
\usepackage{csvsimple}
\usepackage{wrapfig}
\usepackage{comment}
\usepackage{cprotect}
\usepackage{tikz}
\usetikzlibrary{intersections, calc, arrows.meta}
\usetikzlibrary{positioning, through}
\usepackage{fancybox}
\usepackage{booktabs}
\usepackage{adjustbox}
\usepackage{longtable}
\usepackage{nameref}
% 数式系パッケージ
\usepackage{amsmath,amssymb,amsfonts}
\usepackage{mathtools,amssymb,amsfonts}
\usepackage{amsthm}
\usepackage{bm}
% LuaTeX-ja設定
\usepackage{luatexja}
\usepackage[no-math]{luatexja-fontspec}

\setcounter{section}{0}
\makeatletter
    \renewcommand{\theequation}{
    \thesection.\arabic{equation}}
    \@addtoreset{equation}{section}
    \renewcommand{\figurename}{Figure }
    \renewcommand{\tablename}{Table }
\makeatother
\geometry{top=15truemm,bottom=25.4truemm,inner=19.05truemm,outer=19.05truemm}

% Remove section number from equation number, so that it is just 1, 2, NOT 1.1, 2.1:
% \counterwithout{equation}{section}

% Source code listing settings (listings package)
\lstset{
    frame=tb,
    basicstyle=\ttfamily\small,
    % basicstyle=\ttfamily\small\bfseries,
    columns=fullflexible,
    keepspaces=true,
    numbers=left,
    firstnumber=1,
    numberstyle=\scriptsize,
    breaklines=true,
    % frame=none
    % frame=single,
    % gobble=4,
    % xleftmargin=0.2\textwidth,
    % xrightmargin=0.2\textwidth,
}
% \renewcommand{\lstlistlistingname}{List of Source Codes}
% \renewcommand{\lstlistingname}{リスト}

\NewDocumentCommand{\bissectrice}{%
    O{}     % drawing options
    mmm     % bissector of mmm
    m       % intersection point between base and bissector
    O{1}O{1}% extended drawing of the bissector
    }{%
    \path[name path=Bis#2#3#4] let
        \p1 = ($(#2) - (#3)$),          % chktex 46
        \p2 = ($(#4) - (#3)$),          % chktex 46
        \n1 = {veclen(\x1,\y1)},        % chktex 36
        \n2 = {veclen(\x2,\y2)},        % chktex 36
        \n3 = {max(\n1,\n2)},           % chktex 36
        \p1 = ($(#3)!\n3!(#2)$),        % chktex 46
        \p2 = ($(#3)!\n3!(#4)$),        % chktex 46
        \p3 = ($(\p1) + (\p2) - (#3)$)  % chktex 46
    in
        (#3) --(\p3);                   % chktex 36

    \path[name path = foo] (#2) --(#4); % chktex 36

    \path[name intersections={of=foo and Bis#2#3#4, by=#5}];

    \path[#1] ($(#3)!#6!(#5)$) --($(#5)!#7!(#3)$); % chktex 46, chktex 36
}


\counterwithout{equation}{section}
\geometry{top=20truemm,bottom=25.4truemm,inner=19.05truemm,outer=19.05truemm}
\title{\bfseries {\huge Testing My Preamble for Lua{\LaTeX{}}}}
\author{{\LARGE Ivan Setiawan}}
\date{2026年3月28日}

\begin{document}
\maketitle
\hrule
\lstlistoflistings{}
\vspace{0.5em}
\hrule
\vspace{1em}

\section{00preambleの使い方}
\end{document}
\end{lstlisting}

\cprotect\section{The \verb|\counterwithout{}| command}

By default, the equation numbering in \LaTeX{} is like (1.1), (1.2), etc., where the first number is the section number and the second number is the equation number within that section.

If you want to number equations without the section number, use the \verb|\counterwithout{}| command as shown in Listing~\ref{lst:equationnumbernosection}. After using this command, the equation numbering will be like (1), (2), etc., without the section number. Just comment out the \verb|\counterwithout{}| command if you want to revert to the default numbering style, or do not call it in the first place.

\begin{lstlisting}[
    language={[LaTeX] TeX},
    caption={Equation numbering without section number},
    label={lst:equationnumbernosection},
    firstnumber=5,
]
\counterwithout{equation}{section}
\end{lstlisting}

\cprotect\section{The \verb|\geometry{}| command}

To change the size of the page margins, use the \verb|\geometry{}| command from the \texttt{geometry} package. The example in Listing~\ref{lst:geometrycommand} sets the top margin to 20mm, the bottom margin to 25.4mm (1 inch), and both the inner and outer margins to 19.05mm (0.75 inch). Adjust these values as needed for your document.

\begin{lstlisting}[
    language={[LaTeX] TeX},
    caption={\texttt{\textbackslash{}geometry{}} command},
    label={lst:geometrycommand},
    firstnumber=6,
]
\geometry{top=20truemm,bottom=25.4truemm,inner=19.05truemm,outer=19.05truemm}
\end{lstlisting}

You can modify the \verb|\geometry{}| command in your document to set custom margins, or simply comment it out to use the default margins defined in the preamble (file \texttt{00preamble.tex}).

\section{Miscellaneous}

\cprotect\subsection{Use braces \verb|{}| around the \LaTeX{} command to protect it from being separated from the preceeding word by \texttt{latexindent} (\LaTeX{} source file formatter)}

During ``Format Document'', to prevent \texttt{latexindent} from inserting space between a word and the subsequent \LaTeX{} command, place the subsequent command in curly braces directly after the word. For example in the source textfile, write \verb|Lua{\LaTeX{}}| instead of \verb|Lua\LaTeX{}|. The former will be treated as a single unit and will not be separated by \texttt{latexindent}, while the latter may be split into ``Lua'' and ``\verb|\LaTeX{}|'' with a space in between.

\cprotect\subsection{Use \verb|\cprotect| to allow verbatim text in section/subsection titles}

To include verbatim text (e.g., code snippets) in section, subsection titles, or captions, use the \verb|\cprotect| command from the \texttt{cprotect} package. This allows you to write verbatim text without causing errors in the title formatting. For example, you can write:
\begin{lstlisting}[
    language={[LaTeX] TeX},
    caption={Allow verbatim text \texttt{\textbackslash{}verb|text|} in section, subsection titles, or captions},
    label={lst:cprotectallowverbatimtextintitles},
    numbers=none,
]
\cprotect\section{This is a section title with verbatim text: \verb|\LaTeX{}|}
\end{lstlisting}

\cprotect\subsection{Use \verb|\textbackslash{}|}

Use the command \verb|\textbackslash{}| to produce a single backslash character in running text. For example, \verb|This is a backslash: \textbackslash{}| will produce ``This is a backslash: \textbackslash{}'' in the compiled document.

\cprotect\subsection{Use \verb|\url{}| to insert clickable URL link}

To include clickable URL link, write like below:
\begin{lstlisting}[
    language={[LaTeX] TeX},
    caption={How to insert a clickable URL},
    label={lst:urlinsert},
    numbers=none,
]
\url{https://math-note.xyz/latex/tikz/tikz-line/}
\end{lstlisting}
, which will produce \url{https://math-note.xyz/latex/tikz/tikz-line/}.

\cprotect\subsection{Use \texttt{center, Sbox, minipage, verbatim*, shadowbox, fbox, ovalbox} to display \LaTeX{} code snippet in a box}
\label{sec:displaylatexcodeinabox}

\vspace{1em}
\begin{center}
    \begin{Sbox}
        \begin{minipage}{0.8\textwidth}
            \begin{verbatim*}
                \begin{tikzpicture}
                    \draw 〜〜〜〜;
                \end{tikzpicture}\end{verbatim*}
        \end{minipage}
    \end{Sbox}
    \shadowbox{\TheSbox}
\end{center}

The above \LaTeX{} code snippet inside shadowbox is produced using the following code in Listing~\ref{lst:displaylatexcodeinabox}.

\begin{lstlisting}[
    language={[LaTeX] TeX},
    caption={How to display \LaTeX{} code snippet in a box},
    label={lst:displaylatexcodeinabox},
    xleftmargin=0.2\textwidth,
    xrightmargin=0.2\textwidth,
]
\begin{center}
    \begin{Sbox}
        \begin{minipage}{0.8\textwidth}
            \begin{verbatim*}
                \begin{tikzpicture}
                    \draw 〜〜〜〜;
                \end{tikzpicture}\end{verbatim*}
        \end{minipage}
    \end{Sbox}
    \shadowbox{\TheSbox}
\end{center}
\end{lstlisting}

\section[short]{Better to use \textit{center}-ed \texttt{lstlisting}}

Instead of using \texttt{center, Sbox, minipage, verbatim*, shadowbox, fbox, ovalbox} as shown in Subsection~\ref{sec:displaylatexcodeinabox}, it is better to use the \texttt{lstlisting} environment.

The \texttt{center} environment-like is obtained by setting \texttt{xleftmargin, xrightmargin}. The fbox-like is obtained by setting \texttt{frame=single}. Furthermore, redundant whitespace in front of each line can be removed by setting \texttt{gobble} to the number of whitespace characters to remove. For example, if each line starts with 8 spaces, set \texttt{gobble=8}.

\lstnewenvironment{TeXlstlisting}
{\lstset{
        language=[LaTeX]TeX,
        caption={How to use \textit{center}-ed \texttt{lstlisting}},
        label={lst:centeringlstlisting},
        frame=tb,
        basicstyle=\ttfamily\small,
        columns=fullflexible,
        keepspaces=true,
        numbers=left,
        firstnumber=1,
        numberstyle=\scriptsize,
        breaklines=true,
    }
}
{}

% \begin{noindent}
\begin{TeXlstlisting}
\begin{lstlisting}[
    language={[LaTeX] TeX},
    frame=single,
    gobble=8,
    xleftmargin=0.2\textwidth,
    xrightmargin=0.2\textwidth,
]
        ``Now, an example of inline math and TikZ:\
        \(
        \alpha + \beta = \gamma \,
        \begin{tikzpicture}
            \draw (0,0) -- (2,1);
        \end{tikzpicture}
        \)
        inside a sentence''.
\end{lstlisting}
\end{TeXlstlisting}
% \end{noindent}

The above code Listing~\ref{lst:centeringlstlisting} will produce a centered \texttt{lstlisting} with a single frame and no leading whitespace, which is ideal for displaying \LaTeX{} code snippets, like below:
\begin{lstlisting}[
    language={[LaTeX] TeX},
    frame=single,
    gobble=8,
    xleftmargin=0.2\textwidth,
    xrightmargin=0.2\textwidth,
]
        ``Now, an example of inline math and TikZ:\
        \(
        \alpha + \beta = \gamma \,
        \begin{tikzpicture}
            \draw (0,0) -- (2,1);
        \end{tikzpicture}
        \)
        inside a sentence''.
\end{lstlisting}

\end{document}
